\section{A sharp large sieve for a common orbit profile}
\label{sec:tpc48-large-sieve}

The barrier in \cref{thm:tpc48-unrestricted-barrier} permits each slope
to choose an independent compensating orbit vector.  The projective
layers in \eqref{eq:tpc48-projective-layer} do not have that freedom:
all slopes share the same sampled scalar profile.  This section
quantifies the resulting gain.

\subsection{The progression Gram kernel}

Fix \(a\pmod{H_0}\), \(J\ge1\),
\(\phi\in C_c^\infty(\R)\), and a finite slope set
\(\cM\subset[M_-,M_+]\cap\Z\).  For Hilbert vectors
\(b_m\in\Vv\), define
\begin{equation}
 \cL_{\phi,a,\alpha}(\bm b)
 :=\sum_{\substack{j\in\Z\\j\equiv a\!\!\pmod{H_0}}}
 \phi(j/J)
 \left(\sum_{m\in\cM}b_m\e(mj\alpha)\right)\otimes e_j.
 \label{eq:tpc48-common-profile-operator}
\end{equation}
Put \(K=|\cM|\), \(t=\|\alpha\|_\T\), and
\begin{equation}
 D_J(t;K):=
 \begin{cases}
  \min\left\{K,1+\dfrac1{Jt}\right\},&t>0,\\[4pt]
  K,&t=0.
 \end{cases}
 \label{eq:tpc48-collision-factor}
\end{equation}

Let \(\psi=|\phi|^2\).  The Gram kernel is
\begin{equation}
 G_d(\alpha)
 :=\sum_{j\equiv a\!\!\pmod{H_0}}
 \psi(j/J)\e(dj\alpha),
 \qquad d\in\Z.
 \label{eq:tpc48-Gram-kernel}
\end{equation}

\begin{lemma}[No-wrap progression decay]
\label{lem:tpc48-Gram-decay}
Assume \(t\ge0\) and
\begin{equation}
 H_0(M_+-M_-)t\le\frac14.
 \label{eq:tpc48-large-sieve-no-wrap}
\end{equation}
For every integer \(B\ge2\),
\begin{equation}
 |G_d(\alpha)|
 \le C_{B,H_0,\phi}\frac{J}{H_0}
 (1+J|d|t)^{-B}
 \label{eq:tpc48-Gram-decay}
\end{equation}
uniformly for \(|d|\le M_+-M_-\).
\end{lemma}

\begin{proof}
Progression Poisson summation gives
\begin{equation}
 G_d(\alpha)
 =\frac{J}{H_0}\sum_{q\in\Z}
 \e(aq/H_0)
 \widehat\psi\!\left(
 \frac{J}{H_0}(q-H_0d\alpha)
 \right).
 \label{eq:tpc48-Gram-Poisson}
\end{equation}
Choose the representative of \(\alpha\) with absolute value \(t\).
Under \eqref{eq:tpc48-large-sieve-no-wrap}, the \(q=0\) term is
bounded by the right side of \eqref{eq:tpc48-Gram-decay}.  Every
\(q\ne0\) satisfies \(|q-H_0d\alpha|\ge3/4\).  Apply Schwartz decay at
order \(B+2\) to those terms.  If \(J\ge H_0\), their sum after the
prefactor \(J/H_0\) is \(O_{B,H_0,\phi}(J^{-B-1})\), which is bounded
by the right side of \eqref{eq:tpc48-Gram-decay} because
\(1+J|d|t\le1+J/(4H_0)\).  If \(J<H_0\), the conclusion is absorbed
into the constant.  Thus only a fixed finite Schwartz seminorm of
\(\psi\) is used.
\end{proof}

\subsection{The Hilbert-valued inequality}

\begin{theorem}[Common-profile determinant large sieve]
\label{thm:tpc48-common-profile-large-sieve}
Under \eqref{eq:tpc48-large-sieve-no-wrap},
\begin{equation}
 \boxed{
 \|\cL_{\phi,a,\alpha}(\bm b)\|^2
 \le C_{\phi,H_0}\frac{J}{H_0}
 D_J(t;K)\sum_{m\in\cM}\|b_m\|^2.}
 \label{eq:tpc48-common-profile-large-sieve}
\end{equation}
The constant is uniform in the Hilbert space, the vectors \(b_m\),
the slope locations inside \([M_-,M_+]\), and \(J\).
\end{theorem}

\begin{proof}
If \(t=0\), Cauchy--Schwarz in the \(m\)-sum and
\(\sum_{j\equiv a\ (H_0)}|\phi(j/J)|^2\ll_{\phi,H_0}J/H_0\)
give \eqref{eq:tpc48-common-profile-large-sieve} with
\(D_J(0;K)=K\).  Assume henceforth that \(t>0\).
Expanding the square in \eqref{eq:tpc48-common-profile-operator} gives
the Hilbert-valued Gram form
\begin{equation}
 \sum_{m,n\in\cM}
 G_{m-n}(\alpha)\langle b_m,b_n\rangle.
 \label{eq:tpc48-Gram-form}
\end{equation}
By \cref{lem:tpc48-Gram-decay} with \(B=2\),
\begin{equation}
 \sup_m\sum_{n\in\cM}|G_{m-n}(\alpha)|
 \ll_{\phi,H_0}\frac{J}{H_0}
 \left(1+\frac1{Jt}\right).
 \label{eq:tpc48-Gram-row-sum}
\end{equation}
The Schur test, valid after scalar pairing in \(\Vv\), gives this
constant.  The trivial estimate
\(|G_d|\ll_{\phi,H_0}J/H_0\) gives instead
\((J/H_0)K\).  Taking the smaller bound proves
\eqref{eq:tpc48-common-profile-large-sieve}.
\end{proof}

By orbit Plancherel, the left side is also the complete all-band square
function
\begin{equation}
 \|\cL_{\phi,a,\alpha}(\bm b)\|^2
 =\sum_{k\bmod L}\|\Pi_k
 \cL_{\phi,a,\alpha}(\bm b)\|^2
 \label{eq:tpc48-large-sieve-all-band}
\end{equation}
for every choice of \(L\).  In particular, the estimate already
controls the reassembled orbit field and carries no factor equal to the
number of bands.

\subsection{Sharpness in the common-profile class}

\begin{proposition}[Clustered-slope lower bound]
\label{prop:tpc48-large-sieve-sharpness}
Assume \(t>0\), \(\phi\not\equiv0\), and
\(J\ge J_0(\phi,H_0)\).  There are constants
\(c_{\phi,H_0},C_{\phi,H_0}>0\) with the following property.  Suppose
\(\cM\) contains \(D\) consecutive integers and either \(D=1\), or
\begin{equation}
 2\le D\le c_{\phi,H_0}\min\{K,(Jt)^{-1}\}.
 \label{eq:tpc48-coherent-cluster-size}
\end{equation}
Then one can choose collinear Hilbert vectors, supported on those
slopes, such that
\begin{equation}
 \|\cL_{\phi,a,\alpha}(\bm b)\|^2
 \ge C_{\phi,H_0}^{-1}\frac{J}{H_0}D
 \sum_m\|b_m\|^2.
 \label{eq:tpc48-cluster-lower-bound}
\end{equation}
Consequently \(D_J(t;K)\) is optimal up to constants for dense
common-profile slope boxes.
\end{proposition}

\begin{proof}
Choose an interval on which \(|\phi|\) is bounded below.  The lower
bound on \(J\) ensures that its \(J\)-dilate contains
\(\gg_{\phi,H_0}J/H_0\) points in the progression, uniformly in the
finitely many residue classes \(a\).  For \(D=1\), any one supported
slope gives the claim directly.  For \(D\ge2\), let \(j_0\) be the
center of the sampled interval and choose the signed lift
\(\widetilde\alpha\in[-1/2,1/2]\) with
\(|\widetilde\alpha|=t\).  Fix a unit vector \(v\), and on a
consecutive block put
\(b_m=\e(-mj_0\widetilde\alpha)v\).  For \(j\) in a sufficiently
short fixed subinterval around \(j_0\), condition
\eqref{eq:tpc48-coherent-cluster-size} makes all phases
\(\e(m(j-j_0)\widetilde\alpha)\) lie in one fixed acute arc.  Their sum has
modulus \(\gg D\).  Summing over those \(j\) gives
\(\gg_{\phi,H_0}(J/H_0)D^2\), whereas
\(\sum_m\|b_m\|^2=D\).
\end{proof}

Taking \(D=1\) when \(Jt\gg1\), and
\(D\asymp_{\phi,H_0}\min\{K,(Jt)^{-1}\}\) when \(Jt\ll1\), proves the
stated optimality for slope boxes containing the required consecutive
cluster.  At \(t=0\), choosing all \(K\) vectors equal gives the exact
factor \(K\).  The sharpness witness uses arbitrary collinear coefficients
inside one admissible common-profile class.  It does not assert that
the literal actual coefficients realize this coherence.
